% Author: Adam | Module: AI Job Interview Simulator | Project: SmartCareer AI
\section{Performances}
\label{sec:perf}

Toutes les chiffres de cette section viennent de mesures réelles sur la machine
de référence (RTX 3050 Ti Laptop, 4096~Mio de VRAM, 13,9~Go de RAM, modèle
\texttt{mistral:7b-instruct-v0.2-q4\_K\_M} servi par Ollama 0.40.11). Le script
de mesure (\texttt{\_perf\_live.py}, hors dépôt) interroge Ollama en direct avec
le même assemblage de payload que \file{llm\_adam.py}, et récupère les compteurs
renvoyés par le moteur (\texttt{prompt\_eval\_count}, \texttt{eval\_count},
\texttt{eval\_duration}) ; l'échantillonnage mémoire est fait par
\texttt{nvidia-smi} juste après chaque appel.

\subsection{Coût par étape}
\label{sec:perf-latence}

\begin{small}
\begin{table}[H]
\centering
\caption{Les trois étapes du pipeline mesurées en \texttt{num\_gpu=0} (CPU pur),
schéma envoyé uniquement par \texttt{format=}.}
\label{tab:etapes-cpu}
\begin{tabular}{@{}l r r r l r r@{}}
\toprule
\textbf{Étape} & \textbf{Temps} & \textbf{Tok/s} & \textbf{Sortie} & \textbf{Cap} & \textbf{Tok. prompt} & \textbf{Car. prompt} \\
\midrule
Questions ($n=3$) & 113,6 s & 6,3 & 302 t & 630 & 1\,797 & 5\,968 \\
Évaluation        & 94,3 s  & 6,3 & 220 t & 420 & 1\,792 & 5\,425 \\
Amélioration      & 121,5 s & 6,1 & 324 t & 520 & 2\,029 & 6\,591 \\
\bottomrule
\end{tabular}
\end{table}
\end{small}

\textbf{Levier principal : \texttt{num\_predict}, pas la taille du prompt.}
Le préfill mesuré est de 9\,000 à 12\,300~tok/s selon la configuration : un
prompt de 1800~tokens coûte donc $\approx$0,16~s, tandis que 300~tokens de
sortie en CPU en coûtent $\approx$48. Réduire le prompt de 30\,\%
(§\ref{sec:prompts-taille}) ne change la latence que de quelques dixièmes de
seconde ; ce sont les caps par étape --- et la consigne de longueur dans le
prompt --- qui ont fait baisser le temps perçu. Cette mesure a invalidé une
hypothèse de travail initiale, et elle est la raison pour laquelle la
réduction du prompt a été arrêtée là plutôt que poussée au détriment des
consignes.

\subsection{Balayage de \texttt{num\_gpu}}
\label{sec:perf-numgpu}

Même prompt d'évaluation (5425~caractères), modèle rechargé à chaque
configuration après libération, donc chaque ligne part d'un état propre.

\begin{small}
\begin{table}[H]
\centering
\caption{Balayage \texttt{num\_gpu} $\in \{0, 10, 20, -1\}$ (valeurs de
\texttt{ADAM\_MISTRAL\_NUM\_GPU} ; $-1$ = Ollama décide).}
\label{tab:numgpu}
\begin{tabular}{@{}l r r r r r r@{}}
\toprule
\textbf{\texttt{num\_gpu}} & \textbf{Froid} & \textbf{dont charge} & \textbf{Chaud} & \textbf{Gen. tok/s} & \textbf{Préfill tok/s} & \textbf{VRAM} \\
\midrule
0 (CPU pur)   & 110,0 s & 10,9 s & 36,7 s & 6,4  & 11\,073 & \multicolumn{1}{c}{0 MiB} \\
10 couches    & 32,7 s  & 11,9 s & 19,0 s & 7,1  & 9\,424  & 1\,789 MiB \\
20 couches    & 28,0 s  & 8,4 s  & 25,3 s & 9,6  & 12\,261 & \textbf{3\,321 MiB} \\
$-1$ (auto)   & 39,5 s  & 7,6 s  & 21,2 s & 7,4  & 9\,050  & 2\,251 MiB \\
\bottomrule
\end{tabular}
\end{table}
\end{small}

Les quatre configurations ont produit un JSON valide au premier essai. La
lecture des données n'est pas triviale, parce que le critère n'est pas la seule
latence :

\begin{itemize}[itemsep=2pt]
  \item Le GPU accélère réellement : 19--25~s en chaud contre 36,7~s en CPU, et
        le débit de génération monte à 9,6~tok/s avec 20 couches déléguées.
  \item Mais 20 couches consomment \textbf{3321~Mio sur 4096} : il reste
        $\approx$775~Mio, et le Module 1 a besoin du GPU pour Qwen3-0.6B + LoRA.
        Les deux modèles ne peuvent pas coexister dans cet état.
  \item \texttt{num\_gpu=10} est le meilleur compromis mesuré en latence chaude
        (19,0~s) pour une empreinte de 1789~Mio, soit un peu de marge.
  \item \texttt{num\_gpu=-1} (le défaut d'Ollama, donc le défaut du module)
        donne 21,2~s et 2251~Mio : rapide, mais Ollama peut en déléguer davantage
        lorsqu'il croit avoir de la place, ce qui est précisément le scénario
        d'OOM quand le Module 1 charge.
\end{itemize}

\textbf{Recommandation issue de la mesure} (écrite dans \file{.env.example}) :
pour une démonstration ou une soutenance où le Module 1 est \emph{oisif},
garder \texttt{-1} ou \texttt{10} et libérer en fin de séance ; pour un usage où
les deux modules peuvent tourner en concurrence, \texttt{ADAM\_MISTRAL\_NUM\_GPU=0}
--- plus lent d'environ 1,7$\times$ en chaud, mais jamais d'OOM, et l'entretien
ne fait jamais planer le poste. Le défaut du code reste \texttt{-1} : le module
ne présume pas de l'usage.

\subsection{Froid contre chaud, et le plafond de timeout}
\label{sec:perf-timeout}

Sur un même prompt en \texttt{num\_gpu=0} : appel froid (modèle à charger)
\textbf{108,6~s} dont $\approx$10,6~s de chargement pur ; appel chaud
\textbf{37,1~s}. Soit un facteur 2,9 entre le premier appel d'une séance et les
suivants --- et c'est le premier appel qui suit le bouton « Démarrer », donc
celui que l'utilisateur juge.

Le plafond de lecture mérite une explication à part, parce qu'il a produit un
comportement contre-intuitif. \texttt{TIMEOUT\_READ} d'\texttt{requests} est un
délai \emph{maximum entre deux octets}, pas une durée totale --- mais avec
\texttt{stream=false}, Ollama n'écrit la réponse qu'une fois la génération
terminée. Un read timeout vaut donc, ici, plafond de génération. Un appel de
génération de questions avait duré \textbf{272~s} en apparence --- contre 110~s
attendus --- et un seul appel froid bien calibré coûte 109~s ; l'écart s'est
expliqué par la boucle de retry, puis reproduit par la mesure :

\begin{quote}\small
plafond initial 180~s $\to$ tentative 1 froide interrompue à 180~s, backoff
$\approx$3~s, tentative 2 \textbf{chaude} $\approx$89~s $\to$ succès,
\textbf{total $\approx$272~s}. Sur le prompt court de la campagne de mesure
(appel froid 108,6~s, chaud 37,1~s), le même scénario donne $\approx$220~s.
\end{quote}

Le retry transformait donc un appel lent en appel deux fois plus long que
nécessaire, avec une erreur intermédiaire invisible. Le correctif est le plafond
\texttt{ADAM\_TIMEOUT\_READ=600} (l'ancien nom \texttt{ADAM\_MISTRAL\_TIMEOUT\_READ}
reste accepté), qui laisse un appel froid aboutir au lieu de repartir en
retry ; les tentatives supplémentaires (\texttt{ADAM\_MISTRAL\_RETRIES=2},
backoff \texttt{3,0~s}$\times$essai + jitter aléatoire, \file{llm\_adam.py}:175--177)
restent pour les pannes réseau et les 5xx. \texttt{SEMAPHORE\_TIMEOUT} vaut
\texttt{TIMEOUT\_READ + 60}, pour qu'une requête en attente du sémaphore n'abandonne
jamais avant celle qui tient le modèle.

\subsection{Coexistence avec le Module 1 sur 4~Gio}
\label{sec:perf-vram}

Le module ne touche jamais le GPU directement : il ne demande de la VRAM
qu'indirectement, par Ollama. Trois faits mesurés :

\begin{enumerate}[label=\arabic*.,leftmargin=*,itemsep=2pt]
  \item Mistral chargé occupe \textbf{2251~Mio} en configuration automatique
        (2239~Mio relevés lors du parcours navigateur, même ordre).
  \item \texttt{POST /interview/release} (méthode \texttt{release()},
        \file{llm\_adam.py}:215--229, POST avec \texttt{keep\_alive: 0}) aboutit :
        \texttt{ollama ps} revient vide et le GPU est mesuré à
        \textbf{0~Mio / 4096~Mio}. Vérifié depuis le navigateur sur un vrai
        rechargement, sur une navigation SPA vers l'éditeur de CV, et sur une
        fermeture d'onglet.
  \item \texttt{keep\_alive=5m} par défaut fait la même chose sans intervention,
        avec cinq minutes de latence.
\end{enumerate}

\textbf{Limite mesurée, non corrigée.} Une libération demandée
\emph{pendant} une génération en cours ne libère rien : la requête
\texttt{keep\_alive: 0} est acceptée, mais le modèle reste chargé jusqu'à la fin
de la génération en vol, et \texttt{ollama ps} n'est pas vide dans l'intervalle.
C'est ce point qui a été noté en réserve à l'étape~7 du parcours navigateur.
Le module pourrait le forcer (Ollama accepte une annulation de requête), mais
cela signifierait couper une réponse en cours de rédaction --- ce qui a été
jugé préférable de ne pas faire sans le demander à l'utilisateur. Le cas est
documenté en §\ref{sec:limites} comme la principale réserve du volet mémoire.

\subsection{Mesure reproductible pendant la rédaction de ce rapport}
\label{sec:perf-controle}

Quatre appels d'évaluation relancés avec le prompt de contrôle (script
\texttt{calibrate\_eval.py}, \texttt{num\_gpu} par défaut) : 17,7~s ; 21,8~s ;
28,4~s ; 34,1~s, pour des prompts de 4672 et 5182~caractères. Quatre réponses
sur quatre valides au premier essai. Ces valeurs, plus basses que le
\texttt{num\_gpu=0} du tableau~\ref{tab:etapes-cpu}, confirment la déléga-
tion GPU effective dans la configuration par défaut. Le débit effectif dépend
de la longueur réelle de la sortie (bornée à 420~tokens par
\texttt{ADAM\_EVALUATE\_MAX\_TOKENS}), pas seulement de la configuration.

\subsection{Ce qui n'est pas mesuré}
\label{sec:perf-nonsources}

Par honnêteté, les mesures absentes de ce chapitre : aucun test de charge
concourante multi-utilisateurs (\texttt{MAX\_CONCURRENT\_LLM\_CALLS} protège le
modèle, mais son comportement sous 10 utilisateurs simultanés n'a pas été
mesuré) ; aucune latence de la voie d'erreur hors des deux scénarios testés
(Ollama coupé, backend coupé) ; aucun pourcentage de couverture de code
(\texttt{pytest --cov} n'a pas été lancé, cf. §\ref{sec:tests}) ; aucune mesure
de qualité du score sur un jeu annoté (la calibration
§\ref{sec:tests-calibration} porte sur 4 appels, pas sur un jeu d'or) ; aucun
coût énergétique ni température.
